% Compile: pdflatex report -> biber report -> pdflatex report -> pdflatex report
%      or: latexmk -pdf report.tex

\documentclass[11pt,a4paper]{article}

\usepackage[margin=2.5cm]{geometry}
\usepackage{booktabs}
\usepackage[hidelinks]{hyperref}
\usepackage[style=apa,backend=biber]{biblatex}
\addbibresource{references.bib}

\title{Read-Only API Service: Local Read Model vs.\ Per-Request Aggregation}
\author{Your Name \\ Team, Company}
\date{\today}

\begin{document}

\maketitle

\begin{abstract}
\noindent
% Executive summary. One paragraph: question, approach, finding, recommendation.
% Write this last. State the recommendation explicitly.

This is an abstract  \parencite{richardson2018}

\end{abstract}

\tableofcontents

\section{Introduction}

\subsection{Background}
% Why this question came up, and what system it concerns.

\subsection{Problem statement}
% Name the options in the vocabulary the literature uses, so citations follow:
% (i) local read model populated by upstream services (CQRS read model);
% (ii) per-request aggregation of upstream APIs (API composition).

\subsection{Scope}
% In scope / out of scope.

\section{Methodology}

\subsection{Literature search}
% Sources searched, search terms, date range, inclusion and exclusion criteria.

\subsection{Internal inputs}
% Interviews, monitoring data, existing design documents.

\subsection{Evaluation criteria}
% Define criteria here, before presenting findings.

\begin{table}[h]
\centering
\caption{Evaluation criteria.}
\begin{tabular}{@{}lll@{}}
\toprule
Criterion & Definition / how measured & Weight \\
\midrule
Data freshness        & & \\
Read latency          & & \\
Availability          & & \\
Query capability      & & \\
Operational cost      & & \\
Implementation effort & & \\
\bottomrule
\end{tabular}
\end{table}

\subsection{Experimental setup}
% Only if you benchmarked: dataset, load profile, tool, runs, what was constant.

\section{Findings}
% Organise by option or by criterion, never source by source.

\subsection{Option A}

\subsection{Option B}

\subsection{Disagreement between sources}

\subsection{Measurement results}

\section{Analysis}

\subsection{Comparison against criteria}

\begin{table}[h]
\centering
\caption{Option comparison.}
\begin{tabular}{@{}lll@{}}
\toprule
Criterion & Option A & Option B \\
\midrule
Data freshness        & & \\
Read latency          & & \\
Availability          & & \\
Query capability      & & \\
Operational cost      & & \\
Implementation effort & & \\
\bottomrule
\end{tabular}
\end{table}

\subsection{Trade-offs}

\subsection{Risks}

\subsection{Conditions that would change the answer}

\section{Limitations and Assumptions}

\section{Recommendation and Next Steps}

\subsection{Recommendation}

\subsection{Rejected alternatives}

\subsection{Next steps}

% Cite with \cite{key}, e.g. \cite{richardson2018}. Keys are in references.bib.
% Use \nocite{*} only while drafting, to check the bibliography renders.
\printbibliography[heading=bibnumbered]

\appendix

\section{Research Log}
% Every source screened, including rejected ones.

\begin{table}[h]
\centering
\begin{tabular}{@{}llll@{}}
\toprule
Date & Source & Type & Claim taken / reason rejected \\
\midrule
 & & & \\
\bottomrule
\end{tabular}
\end{table}

\section{Raw Data}

\end{document}
